\begin{tikzpicture}[x=1cm,y=-.87cm,every node/.style={font=\fontsize{12.2}{15}\selectfont,text=Ink},arr/.style={-{Stealth[length=2mm]},line width=.8pt,draw=Classical}]\path[use as bounding box] (0,0) rectangle (7.45,6.2);
\draw[draw=Classical,fill=Classical!8,rounded corners=2pt,line width=.7pt] (0.25,0.3) rectangle (7.15,2.8);
\node[anchor=center,text=Classical,font=\fontsize{12.2}{15}\selectfont,text width=6.5cm,align=center] at (3.7,0.75) {\textbf{Lista Python}};
\node[anchor=center,text=Ink,font=\fontsize{13}{15}\selectfont,text width=2.5cm,align=center] at (1.85,1.55) {[1, 2]};
\draw[arr,draw=Classical] (3.1,1.55)--(3.8,1.55);
\node[anchor=center,text=Ink,font=\fontsize{13}{15}\selectfont,text width=3cm,align=center] at (5.45,1.55) {[1, 2, 1, 2]};
\node[anchor=center,text=Ink,font=\fontsize{11.5}{15}\selectfont,text width=6.5cm,align=center] at (3.7,2.35) {× 2 ripete gli elementi};
\draw[draw=Quantum,fill=Quantum!8,rounded corners=2pt,line width=.7pt] (0.25,3.25) rectangle (7.15,5.75);
\node[anchor=center,text=Quantum,font=\fontsize{12.2}{15}\selectfont,text width=6.5cm,align=center] at (3.7,3.7) {\textbf{Array NumPy}};
\node[anchor=center,text=Ink,font=\fontsize{13}{15}\selectfont,text width=2.5cm,align=center] at (1.85,4.5) {[1, 2]};
\draw[arr,draw=Quantum] (3.1,4.5)--(3.8,4.5);
\node[anchor=center,text=Ink,font=\fontsize{13}{15}\selectfont,text width=3cm,align=center] at (5.45,4.5) {[2, 4]};
\node[anchor=center,text=Ink,font=\fontsize{11.5}{15}\selectfont,text width=6.5cm,align=center] at (3.7,5.3) {× 2 raddoppia ogni valore};

\end{tikzpicture}